\section{Periodic density fitting and RSDF response}
\label{sec:periodic-df}

\subsection{Why density-fitting derivatives are required}

Direct treatment of four-center electron-repulsion integrals causes the cost
and memory of periodic HF and target-DFA Coulomb and exact-exchange energies to
grow rapidly.  Density fitting expands AO-pair densities in an auxiliary
Gaussian basis and replaces the four-center problem by three-center quantities
and a two-center metric.  When the nuclei move, however, both the primary AOs
and the auxiliary basis move with them.  Thus, if density fitting is used for
the energy value, the gradient must include the auxiliary-basis and metric
responses.

Denote a Bloch AO by $\chi_{\mu\vk}(\mathbf r)$ and define the momentum transfer
as
\begin{equation}
  \mathbf q=\mathbf k_j-\mathbf k_i.
\end{equation}
For an auxiliary Bloch function $P\mathbf q$, the three-center source and the
Coulomb metric are
\begin{align}
  B_{\mu\nu,P}^{ij}
  &=\left(\chi_{\mu\mathbf k_i}\chi_{\nu\mathbf k_j}\middle|P\mathbf q\right),\\
  M_{PQ}(\mathbf q)
  &=\left(P\mathbf q\middle|Q\mathbf q\right).
\end{align}
Whether the ket AO is represented with an explicit complex conjugate depends
on the integral-orientation convention of the program.  The essential
requirement is that the central value and its derivative use the same
convention.

For a positive-definite metric, let
\begin{equation}
  \mathbf M(\mathbf q)=\mathbf L(\mathbf q)\mathbf L(\mathbf q)^\dagger
\end{equation}
be its Cholesky decomposition.  The orthonormalized three-center factor is
\begin{equation}
  \mathbf C^{ij}=\mathbf L(\mathbf q)^{-1}\mathbf B^{ij}.
  \label{eq:cderi-definition}
\end{equation}
A four-center integral is then approximated by a factorization of the form
\begin{equation}
  (\mu i,\nu j|\lambda l,\sigma m)
  \simeq \sum_P C_{P,\mu\nu}^{ij}
  C_{P,\lambda\sigma}^{lm*}.
\end{equation}
The factor $\mathbf C$ corresponds to PySCF's CDERI.

\subsection{Nuclear derivatives and auxiliary response}

Differentiating Eq.~\eqref{eq:cderi-definition} with respect to a Cartesian
coordinate $R_{Ax}$ gives
\begin{equation}
  \mathbf C^{Ax}
  =\mathbf L^{-1}\mathbf B^{Ax}
  -\mathbf L^{-1}\mathbf L^{Ax}\mathbf C.
  \label{eq:cderi-response}
\end{equation}
The first term is the three-center source response and the second is the
metric response.  Freezing the metric omits the motion of the auxiliary
Gaussian centers and the change of the fitting coefficients.  This omission
may be accidentally small for a small numerator, but it generally grows as
the number of auxiliary centers increases.

The Cholesky derivative satisfies
\begin{equation}
  \mathbf M^{Ax}=\mathbf L^{Ax}\mathbf L^\dagger
  +\mathbf L\mathbf L^{Ax\dagger}.
  \label{eq:cholesky-response}
\end{equation}
Multiplying by $\mathbf L^{-1}$ and $\mathbf L^{-\dagger}$ defines
\begin{equation}
  \mathbf W=\mathbf L^{-1}\mathbf M^{Ax}\mathbf L^{-\dagger}.
\end{equation}
In the lower-Cholesky gauge,
\begin{equation}
  \mathbf L^{Ax}=\mathbf L\,\Phi(\mathbf W),
\end{equation}
where $\Phi$ preserves the strict lower triangle and takes half of the
Hermitian diagonal.  The same gauge must be used consistently for the central
value and the derivative.

\subsection{Range separation of the Coulomb kernel}

RSDF partitions the Coulomb kernel according to
\begin{equation}
  \frac{1}{r}
  =\frac{\operatorname{erfc}(\omega r)}{r}
  +\frac{\operatorname{erf}(\omega r)}{r}.
  \label{eq:range-separation}
\end{equation}
The first term decays rapidly in real space and is evaluated as a lattice sum;
the second is evaluated as a smooth reciprocal-space Fourier sum.  This is an
exact identity, but with finite image lists and reciprocal meshes the value of
$\omega$, the screening precision, and the mesh become part of the numerical
representation.

In the public PySCF RSDF implementation, optimal parameters for the
three-center source and the two-center metric may be determined independently.
Explicitly,
\begin{align}
  \mathbf B^{ij}
  &=\mathbf B_{\mathrm{SR}}^{ij}(\omega_B,\epsilon_R)
  +\mathbf B_{\mathrm{LR}}^{ij}(\omega_B,\mathcal G_B)
  +\delta_{\mathbf q,0}\mathbf B_{vbar}^{ij},
  \label{eq:rsdf-source}\\
  \mathbf M(\mathbf q)
  &=\mathbf M_{\mathrm{SR}}(\mathbf q;\omega_M,\epsilon_M)
  +\mathbf M_{\mathrm{LR}}(\mathbf q;\omega_M,\mathcal G_M)
  +\delta_{\mathbf q,0}\mathbf M_{G=0}.
  \label{eq:rsdf-metric}
\end{align}
In general, $\omega_B\neq\omega_M$ and
$\mathcal G_B\neq\mathcal G_M$.  The exponent distribution of the auxiliary
basis differs from that of the primary AO pairs, and the metric may therefore
require a higher cutoff.

At fixed cell, lattice and reciprocal vectors do not change under nuclear
motion.  Holding $\omega$ and the mesh fixed as numerical inputs gives
\begin{align}
  \mathbf B^{ij,Ax}
  &=\mathbf B_{\mathrm{SR}}^{ij,Ax}
  +\mathbf B_{\mathrm{LR}}^{ij,Ax}
  +\delta_{\mathbf q,0}\mathbf B_{vbar}^{ij,Ax},\\
  \mathbf M^{Ax}
  &=\mathbf M_{\mathrm{SR}}^{Ax}
  +\mathbf M_{\mathrm{LR}}^{Ax}
  +\delta_{\mathbf q,0}\mathbf M_{G=0}^{Ax}.
\end{align}
For stress, the lattice and reciprocal mesh themselves vary; the fixed-cell
formulas in this section are therefore not sufficient for a cell derivative.

\subsection{Ownership of primary-AO and auxiliary centers}

A three-center integral contains three atom-centered objects.  Symbolically,
\begin{equation}
  B_{\mu\nu,P}^{Ax}
  =\delta_{A,A_\mu}\partial_{A_\mu x}B_{\mu\nu,P}
  +\delta_{A,A_\nu}\partial_{A_\nu x}B_{\mu\nu,P}
  +\delta_{A,A_P}\partial_{A_Px}B_{\mu\nu,P}.
\end{equation}
Integral-library operators such as \texttt{ip1} and \texttt{ip2} often return
derivatives with respect to an electronic coordinate.  Translational
identities then introduce the opposite sign for a derivative that moves a
nuclear center.  The three contributions must be accumulated after applying
the atom slices for the AO row, AO column, and auxiliary index.

The two-center metric has two auxiliary centers,
\begin{equation}
  M_{PQ}^{Ax}
  =\delta_{A,A_P}\partial_{A_Px}M_{PQ}
  +\delta_{A,A_Q}\partial_{A_Qx}M_{PQ}.
\end{equation}
The short-range metric derivative must use the same screened lattice-image list
as the central value.  Different image cutoffs may both approximate the
continuous Coulomb metric, but they do not define an exact derivative of the
same finite calculation.

\subsection{The \texorpdfstring{$G=0$}{G=0} mode and charge corrections}

The $G=0$ mode of the periodic Coulomb kernel is treated specially according
to charge neutrality and the Ewald convention.  The source \texttt{vbar} and
the metric $G=0$ correction occur only in blocks with exactly zero momentum
transfer,
\begin{equation}
  \delta_{\mathbf q,0}=\begin{cases}
    1,& \mathbf k_j-\mathbf k_i=\mathbf 0,\\
    0,& \text{otherwise}.
  \end{cases}
\end{equation}
This includes not only Gamma-only calculations but also diagonal pairs in a
sampled-k calculation.  Therefore, ``is this Gamma?'' and ``is the momentum
transfer zero?'' are not equivalent tests.

\subsection{k-point pairs and Hermitian reconstruction}

For $N_k$ k-points, the J/K build requires multiple $(i,j)$ pairs.  Momentum
conservation groups the metric by $\mathbf q=\mathbf k_j-\mathbf k_i$.  If only
canonical pairs are stored, a reverse pair is reconstructed through a
Hermitian relation such as
\begin{equation}
  C^{ji}_{P,\nu\mu}=C^{ij*}_{P,\mu\nu}.
\end{equation}
Its nuclear derivative must obey the same conjugation and AO-index transpose.
At complex k-points, all pair contractions and the equal-weight $1/N_k$
normalization must be consistent for the final gradient to be real.

Define the closed-shell Coulomb auxiliary density by
\begin{equation}
  j_P=\sum_{\vk\mu\nu}C^\vk_{P,\mu\nu}D^\vk_{\nu\mu}.
\end{equation}
Its fixed-density J derivative is
\begin{equation}
  (J^\vk)^{Ax}_{\mu\nu}
  =\frac{1}{N_k}\sum_P
  \left[j_P^{Ax}C^\vk_{P,\mu\nu}
  +j_PC^{\vk,Ax}_{P,\mu\nu}\right].
\end{equation}
Exchange contains two CDERI factors, so the product rule must be applied to
each factor.

\subsection{Role within the DC-DFT gradient}

RSDF determines the numerical representation of the two-electron value and
response; the nonvariational structure of DC-DFT itself is independent of the
backend.  The total gradient may be written as
\begin{equation}
  E_{\mathrm{DC}}^{Ax}
  =E_{\mathrm{target,explicit}}^{Ax}
  +\mathbf z^T\mathbf B_{\mathrm{HF}}^{Ax}.
  \label{eq:dc-rsdf-total}
\end{equation}
The first term contains the target one-electron, Hartree, XC, overlap/Pulay,
and nuclear contributions.  The second includes the orbital response of the HF
density to nuclear motion through the adjoint Z-vector.  RSDF $C^{Ax}$ enters
both the explicit target J derivative and the source HF response operator.

\subsection{Validation principles}

Correctness of the analytic derivative and physical convergence are separate
questions.  First compare with a central finite difference of the same finite
RSDF representation,
\begin{equation}
  g_{Ax}^{\mathrm{FD}}(h)
  =\frac{E(\mathbf R+h\mathbf e_{Ax})
       -E(\mathbf R-h\mathbf e_{Ax})}{2h}.
\end{equation}
Then converge $h$, the source and metric meshes, the primary and auxiliary
bases, and the k-mesh independently.

In addition, a separately reconstructed CDERI must agree with the SCF storage,
\begin{equation}
  \Delta_C=\max_{ijP\mu\nu}
  \left|C_{P,\mu\nu}^{ij,\mathrm{reconstructed}}
  -C_{P,\mu\nu}^{ij,\mathrm{storage}}\right|.
\end{equation}
If $\Delta_C$ is not small, even a mathematically correct derivative formula
is differentiating a representation with a different $\omega$, mesh, pair
orientation, or factor gauge.  Such a derivative must be rejected for
production use.
